% Source: https://artofproblemsolving.com/community/c3h2004310_absolute_value_question
% Stable identity: aops-topic-2004310-q01
\textbf{(a)} By definition, $\sqrt{t}$ (for $t\ge0$) is the unique \emph{nonnegative} number whose square is $t$. The number $|x|^5$ is nonnegative and $(|x|^5)^2=|x|^{10}=x^{10}$, since an even power does not see the sign. Hence $\sqrt{x^{10}}=|x|^5$ $\bigl(=|x^5|\bigr)$.

\textbf{(b)} By (a), $\sqrt{x^{10}}=x^5$ holds exactly when $|x|^5=x^5$, that is, when $x^5\ge0$, i.e.\ $\boxed{x\ge0}$. For example, at $x=-1$ we get $\sqrt{(-1)^{10}}=1$ but $(-1)^5=-1$. So the absolute value matters whenever $x$ may be negative; it can be dropped only when $x$ is known to be nonnegative.
